\section{Module-Level Graph Structures}

Two natural module-level graph structures in \texttt{Mathlib} emerge from the practice of formalizing mathematics in Lean. The first comes from the file system: module names reflect their paths in the repository, forming a rooted tree (\S\ref{sec:module-names}--\ref{sec:namespace-tree}). The second comes from mathematical dependencies: each \texttt{import} statement declares that one module depends on another, forming a directed acyclic graph (\S\ref{sec:import-graph}). The former is a continuation, in the formal setting, of the way mathematicians have long classified their subject---categories such as algebra, analysis, and topology were established well before formalization, and Mathlib's namespace hierarchy faithfully inherits this tradition. The latter captures the actual logical dependencies between files: one module requires definitions and results from another. 

Such dependencies also exist in informal mathematics, analogous to citations between papers, but there they are implicit and incomplete. Formalization makes them explicit, exhaustive, and computable.

\subsection{Module Names}
\label{sec:module-names}

A \emph{module name} in Mathlib is a dot-separated sequence of identifiers, such as
\[
  \texttt{Mathlib.Algebra.Group.Defs}.
\]
Each module name corresponds to a \texttt{.lean} source file, with dots replacing path separators. We write $\mathcal{M}$ for the set of all module names in Mathlib:
\begin{align*}
  \mathcal{M} = \{\; &\texttt{Mathlib.Init},\\
    &\texttt{Mathlib.Algebra.Group.Defs},\\
    &\texttt{Mathlib.Analysis.Normed.Group.Basic},\; \ldots\,\}.
\end{align*}
All statistics in this paper are computed from commit \texttt{534cf0b} (2 February 2026) of the Mathlib repository, using the official \texttt{importGraph} tool~\cite{importgraph}. At this commit, $|\mathcal{M}| = 7{,}563$. The set $\mathcal{M}$ serves as the shared vertex set for both graphs defined in this section.

\begin{definition}[Depth]
The \emph{depth} of a module name is the number of dot-separated components. Formally, if a module name is $c_1.c_2.\cdots.c_k$, then its depth is $k$.
\end{definition}

\begin{example}
\leavevmode
\begin{itemize}
  \item $\texttt{Mathlib}$ has depth~$1$.
  \item $\texttt{Mathlib.Algebra}$ has depth~$2$.
  \item $\texttt{Mathlib.Algebra.Group.Defs}$ has depth~$4$.
\end{itemize}
The maximum depth in Mathlib is~$7$.\footnote{Attained by, e.g., \texttt{Mathlib.CategoryTheory.Limits.Shapes.Pullback.IsPullback.Kernels}.}
\end{example}

We recall the standard graph-theoretic notions used throughout this paper.

\begin{definition}[Directed graph]
A \emph{directed graph} is a pair $G = (V, E)$ where $V$ is a set (called the \emph{vertex set}) and $E \subseteq V \times V$ (called the \emph{edge set}).
\end{definition}

\begin{definition}[Path]
A \emph{path} in a directed graph $G = (V, E)$ is a sequence of vertices $(v_0, v_1, \ldots, v_k)$ such that $(v_i, v_{i+1}) \in E$ for all $0 \le i < k$. The \emph{length} of the path is~$k$.
\end{definition}

\begin{definition}[Cycle]
A \emph{cycle} is a path $(v_0, v_1, \ldots, v_k)$ with $k \ge 1$ and $v_k = v_0$.
\end{definition}

\begin{definition}[Directed acyclic graph]
A directed graph $G = (V, E)$ is a \emph{directed acyclic graph (DAG)} if $G$ contains no cycle.
\end{definition}

\begin{definition}[Rooted tree]
A \emph{rooted tree} is a DAG $T = (V, E)$ with a distinguished vertex $r \in V$ (the \emph{root}) such that for every $v \in V$ there is a unique directed path from $r$ to $v$.
\end{definition}

\subsection{The Namespace Tree}
\label{sec:namespace-tree}

Module names are organized hierarchically: each dot-separated name mirrors a path in the source repository.

\begin{figure}[H]
\centering
\begin{subfigure}[t]{0.42\textwidth}
\centering
\begin{forest}
  for tree={
    font=\scriptsize\ttfamily,
    grow'=0,
    child anchor=west,
    parent anchor=south,
    anchor=west,
    calign=first,
    edge path={
      \noexpand\path [draw, \forestoption{edge}]
        (!u.south west) +(7.5pt,0) |- (.child anchor)\forestoption{edge label};
    },
    before typesetting nodes={
      if n=1 {insert before={[,phantom]}} {}
    },
    fit=band,
    before computing xy={l=13pt},
  }
[Mathlib/
  [Algebra/
    [Group/
      [Defs.lean]
    ]
  ]
  [Data/
    [Nat/
      [Notation.lean]
    ]
  ]
  [\ldots, edge={draw, densely dotted}]
]
\end{forest}
\caption{Directory structure}
\label{fig:dir-tree}
\end{subfigure}%
\hfill
\begin{subfigure}[t]{0.52\textwidth}
\centering
\begin{tikzpicture}[
  every node/.style={font=\scriptsize\ttfamily, inner sep=2pt},
  dot/.style={circle, fill=black, minimum size=3pt, inner sep=0pt},
  >=Stealth, every edge/.style={draw, ->, thin},
]
  \node[dot, label=above:{Mathlib}]              (root) at (1.8,0) {};
  \node[dot, label=left:{Algebra}]               (alg)  at (0,-1.3) {};
  \node[dot, label=right:{Data}]                 (dat)  at (3.6,-1.3) {};
  \node[dot, label=left:{Algebra.Group}]         (grp)  at (0,-2.6) {};
  \node[dot, label=right:{Data.Nat}]             (nat)  at (3.6,-2.6) {};
  \node[dot, label=below:{Algebra.Group.Defs}]   (gd)   at (0,-3.9) {};
  \node[dot, label=below:{Data.Nat.Notation}]    (nn)   at (3.6,-3.9) {};
  \draw (root) -- (alg);
  \draw (root) -- (dat);
  \draw (alg)  -- (grp);
  \draw (dat)  -- (nat);
  \draw (grp)  -- (gd);
  \draw (nat)  -- (nn);
  \node[font=\scriptsize, text=gray] at (1.8,-1.3) {$\cdots$};
\end{tikzpicture}
\caption{Namespace tree $T$: each edge $\in E_T$}
\label{fig:namespace-tree}
\end{subfigure}
\caption{A fragment of the Mathlib source tree (commit \texttt{534cf0b}). Left: the directory layout. Right: the corresponding rooted tree~$T$. Each directory path maps to a module name with dots replacing separators, e.g., \texttt{Mathlib/\allowbreak Algebra/\allowbreak Group/\allowbreak Defs.lean} $\leftrightarrow$ \texttt{Mathlib.\allowbreak Algebra.\allowbreak Group.\allowbreak Defs}.}
\label{fig:prefix-tree}
\end{figure}

We say $m$ is a \emph{prefix} of $m'$ if $m'$ extends $m$ by additional components. For example, $\texttt{Mathlib.Algebra}$ is a prefix of $\texttt{Mathlib.Algebra.Group.Defs}$.

The prefix relation induces a rooted tree $T = (\mathcal{M}, E_T)$ with root $\texttt{Mathlib}$, where
\[
  E_T = \{(m, m') \in \mathcal{M} \times \mathcal{M} \mid m \text{ is a prefix of } m' \text{ and } \mathrm{depth}(m') = \mathrm{depth}(m) + 1\}.
\]

The top-level directories (children of the root) include \texttt{Algebra} (1257 files), \texttt{CategoryTheory} (982 files), \texttt{Analysis} (755 files), etc.

\subsection{The Import Graph}
\label{sec:import-graph}

Each Lean source file begins with \texttt{import} statements that declare its dependencies. For example, the file \texttt{Mathlib/Data/Nat/Notation.lean} contains:

\begin{minted}[frame=single,fontsize=\small,breaklines,bgcolor=gray!5]{lean4}
module

public import Mathlib.Init
\end{minted}

\noindent
This file depends on exactly one other module: \texttt{Mathlib.Init}.

A file with more dependencies is \texttt{Mathlib/Algebra/Group/Defs.lean}, which imports nine modules (shown here in abbreviated form):

\begin{minted}[frame=single,fontsize=\small,breaklines,bgcolor=gray!5]{lean4}
module

public import Batteries.Logic
public import Mathlib.Algebra.Notation.Defs
public import Mathlib.Algebra.Regular.Defs
  -- ... 6 more imports
\end{minted}

\noindent
Each file is a \emph{vertex} in the import graph; each \texttt{import} statement is a directed \emph{edge} from the importing module to the imported module.

Lean~4 supports both \texttt{import} and \texttt{public import}; the latter re-exports the imported module to downstream dependents. For the purposes of the import graph, both forms create an edge---they differ only in visibility, not in the dependency relationship.

For each $m \in \mathcal{M}$, let $\mathrm{imports}(m) \subseteq \mathcal{M}$ denote the set of modules directly imported by $m$ (via either \texttt{import} or \texttt{public import}).

\begin{definition}[Import graph]
The \emph{import graph} of Mathlib is $G_{\mathrm{import}} = (\mathcal{M}, E_{\mathrm{import}})$ where
\[
  E_{\mathrm{import}} = \{(m_1, m_2) \in \mathcal{M} \times \mathcal{M} \mid m_2 \in \mathrm{imports}(m_1)\}.
\]
\end{definition}

Since \texttt{Mathlib/Data/Nat/Notation.lean} contains \texttt{public import Mathlib.Init}, we have the edge $(\texttt{Data.Nat.Notation},\; \texttt{Init}) \in E_{\mathrm{import}}$:

\begin{figure}[H]
\centering
\begin{tikzpicture}[
  every node/.style={font=\small\ttfamily, inner sep=2pt},
  dot/.style={circle, fill=black, minimum size=3pt, inner sep=0pt},
  >=Stealth, every edge/.style={draw, ->, thin},
]
  \node[dot, label=above:{Init}]              (b) at (0,0) {};
  \node[dot, label=above:{Data.Nat.Notation}] (a) at (4,0) {};
  \draw (a) -- (b);
\end{tikzpicture}
\caption{The simplest import edge: \texttt{Data.Nat.Notation} imports \texttt{Init}. All names have the \texttt{Mathlib.}\ prefix removed.}
\label{fig:single-import-edge}
\end{figure}

\begin{remark}
The namespace tree $T$ and the import graph $G_{\mathrm{import}}$ are distinct structures on the same vertex set $\mathcal{M}$. In $T$, edges connect a module to its namespace parent; in $G_{\mathrm{import}}$, edges reflect explicit \texttt{import} dependencies. Import edges need not respect the namespace hierarchy.
\end{remark}

\begin{figure}[ht]
\centering
\begin{tikzpicture}[
  every node/.style={font=\small\ttfamily, inner sep=2pt},
  dot/.style={circle, fill=black, minimum size=3pt, inner sep=0pt},
  >=Stealth, every edge/.style={draw, ->, thin},
]
  \node[dot, label=above:{Init}]                (d) at (0,0) {};
  \node[dot, label=left:{Data.Int.Notation}]    (b) at (-1.5,-1.5) {};
  \node[dot, label=right:{Data.Nat.Notation}]   (c) at (1.5,-1.5) {};
  \node[dot, label=below:{Algebra.Group.Defs}]  (a) at (0,-3) {};

  \draw (a) -- (b);
  \draw (a) -- (c);
  \draw (b) -- (d);
  \draw (c) -- (d);
\end{tikzpicture}
\caption{A fragment of $G_{\mathrm{import}}$: \texttt{Algebra.Group.Defs} imports two modules that both depend on \texttt{Init}, forming a diamond. All names have the \texttt{Mathlib.}\ prefix removed.}
\label{fig:import-subgraph}
\end{figure}

The import graph $G_{\mathrm{import}}$ is acyclic: we verified this by computing a topological ordering of all $7{,}563$ modules using the \texttt{importGraph} tool~\cite{importgraph}.\footnote{The \texttt{importGraph} tool outputs edges in the reverse direction (imported $\to$ importing). Our definition follows the standard dependency convention (importing $\to$ imported); acyclicity is invariant under edge reversal.} In particular, $G_{\mathrm{import}}$ is a DAG.

\subsection{Transitive Reduction and Redundant Imports}
\label{sec:transitive-reduction}

Since $G_{\mathrm{import}}$ is a DAG, it has a unique \emph{transitive reduction}~\cite{ahu1972}: the smallest subgraph $G_{\mathrm{import}}^{-} \subseteq G_{\mathrm{import}}$ with the same reachability relation. An edge $(m_1, m_2) \in E_{\mathrm{import}}$ is \emph{transitively redundant} if $m_2$ is reachable from $m_1$ via a path of length $\ge 2$ in $G_{\mathrm{import}}$---that is, the dependency is already implied by other imports.

\begin{definition}[Transitive reduction]
The \emph{transitive reduction} of $G_{\mathrm{import}}$ is $G_{\mathrm{import}}^{-} = (\mathcal{M}, E_{\mathrm{import}}^{-})$ where
\[
  E_{\mathrm{import}}^{-} = E_{\mathrm{import}} \setminus \{(m_1, m_2)\in E_{\mathrm{import}} \mid \exists\, \text{path } m_1 \to \cdots \to m_2 \text{ in } G_{\mathrm{import}} \text{ of length} \ge 2\}.
\]
\end{definition}

\begin{figure}[H]
\centering
\begin{subfigure}[t]{0.42\textwidth}
\centering
\vspace{0pt}
\begin{tikzpicture}[
  every node/.style={font=\scriptsize\ttfamily, inner sep=1pt},
  dot/.style={circle, fill=black, minimum size=3pt, inner sep=0pt},
  >=Stealth, every edge/.style={draw, ->, thin},
]
  \node[dot, label=above:{Init}]                (d) at (0,0) {};
  \node[dot, label=left:{Data.Int.Notation}]    (b) at (-1.2,-1.2) {};
  \node[dot, label=right:{Data.Nat.Notation}]   (c) at (1.2,-1.2) {};
  \node[dot, label=below:{Algebra.Group.Defs}]  (a) at (0,-2.4) {};

  \draw (a) -- (b);
  \draw (a) -- (c);
  \draw (b) -- (d);
  \draw (c) -- (d);
  \draw[red, dashed, thick] (a) -- (d);
\end{tikzpicture}
\caption{Subgraph of $G_{\mathrm{import}}$}
\label{fig:tr-graph}
\end{subfigure}%
\hfill
\begin{subfigure}[t]{0.46\textwidth}
\vspace{0pt}\raggedleft
\scriptsize\ttfamily
\begin{tabular}[t]{@{}l@{}}
  \textnormal{\sffamily\bfseries Data/Int/Notation.lean} \\[1pt]
  \hspace{1em}public import Init \\[6pt]
  \textnormal{\sffamily\bfseries Data/Nat/Notation.lean} \\[1pt]
  \hspace{1em}public import Init \\[6pt]
  \textnormal{\sffamily\bfseries Algebra/Group/Defs.lean} \\[1pt]
  \hspace{1em}public import Data.Int.Notation \\
  \hspace{1em}public import Data.Nat.Notation \\
  \hspace{1em}\textcolor{red}{public import Init}
    \textnormal{\sffamily\color{red} $\leftarrow$ would be redundant} \\
\end{tabular}
\caption{Source file headers}
\label{fig:tr-source}
\end{subfigure}
\caption{Illustrating transitive redundancy on the diamond from Figure~\ref{fig:import-subgraph}. If \texttt{Algebra.Group.Defs} were to import \texttt{Init} directly (dashed red edge), that edge would be transitively redundant: \texttt{Init} is already reachable via both \texttt{Data.Int.Notation} and \texttt{Data.Nat.Notation}. All names have the \texttt{Mathlib.}\ prefix removed.}
\label{fig:transitive-reduction}
\end{figure}

The raw import graph has $|E_{\mathrm{import}}| = 23{,}570$ edges. After transitive reduction, $|E_{\mathrm{import}}^{-}| = 19{,}448$ edges, removing $4{,}122$ redundant edges ($17.5\%$). These redundant imports are syntactically present in the source files but logically unnecessary: the Lean compiler resolves them through transitive dependencies regardless.

The \texttt{importGraph} tool applies transitive reduction by default when generating its output. We verified \textcolor{red}{(via our own parser)} that its reduced graph ($20{,}881$ edges including the umbrella root \texttt{Mathlib}) agrees exactly with our independently computed reduction on the shared $7{,}563$ modules: all $19{,}448$ edges in $G_{\mathrm{import}}^{-}$ appear in the tool's output, and the $1{,}433$ additional edges in the tool's output are precisely the edges from the umbrella module \texttt{Mathlib} to its direct children.
